% Źródła: Eves s. 287; Wikipedia: Nikolai Lobachevsky, Carl Friedrich Gauss.
% https://en.wikipedia.org/wiki/Nikolai_Lobachevsky
% Uwaga: Wikipedia podaje rok urodzenia 1792, Eves 1793 (s. 287).

\section{Nikołaj Łobaczewski} % lang-pl
\section{Nikolaj Lobachevskij} % lang-it
\label{sekcja_lobaczewski}

Trzeci odkrywca mieszka najdalej od pozostałych i jako pierwszy ogłosi swoje wyniki drukiem. % lang-pl
W tej sekcji zobaczymy, skąd dowie się o problemie, kiedy zabierze głos i dlaczego to jego nazwisko dostanie geometria. % lang-pl
Il terzo scopritore vivrà più lontano degli altri e sarà il primo a rendere pubblici i suoi risultati. % lang-it
In questa sezione vedremo da dove verrà a conoscere il problema, quando prenderà la parola e perché la geometria porterà il suo nome. % lang-it

Nikołaj Iwanowicz Łobaczewski urodzi się w 1792 roku (Eves podaje 1793) niedaleko Niżnego Nowogrodu i zacznie studia w Kazaniu pod kierunkiem Johanna Christiana Martina Bartelsa, przyjaciela Gaussa. % lang-pl
Zostanie profesorem w wieku trzydziestu lat, a od 1827 roku rektorem Uniwersytetu Kazańskiego. % lang-pl
Nikolaj Ivanovič Lobachevskij nascerà nel 1792 (Eves indica il 1793) vicino a Nižnij Novgorod e comincerà gli studi a Kazan' sotto la guida di Johann Christian Martin Bartels, amico di Gauss. % lang-it
Diventerà professore a trent'anni e dal 1827 rettore dell'Università di Kazan'. % lang-it
\index[persons]{Łobaczewski, Nikołaj Iwanowicz} % lang-pl
\index[persons]{Lobachevskij, Nikolaj Ivanovič} % lang-it
\index[persons]{Bartels, Johann Christian Martin}%

Swoje idee przedstawi po raz pierwszy 23 lutego 1826 roku, a wydrukuje je w latach 1829--1830 w publikacjach uniwersytetu kazańskiego, czyli osobno od świata naukowego -- oddzielony od niego, jak zauważy Eves, barierami odległości i języka \cite[s. 287]{eves1_1972}. % lang-pl
Wersję niemiecką, \emph{Geometrische Untersuchungen}, opublikuje w 1840 roku. % lang-pl
Presenterà le sue idee per la prima volta il 23 febbraio 1826 e le stamperà nel 1829--1830 nelle pubblicazioni dell'università di Kazan', cioè separatamente dal mondo scientifico -- diviso da esso, come osserva Eves, da barriere di distanza e di lingua \cite[p. 287]{eves1_1972}. % lang-it
La versione tedesca, \emph{Geometrische Untersuchungen}, la pubblicherà nel 1840. % lang-it

W odróżnieniu od Bolyaia zbuduje wyłącznie geometrię, w której przez punkt przechodzi więcej niż jedna równoległa. % lang-pl
Ponieważ opublikuje ją pierwszy, to jego imię dostanie \emph{geometria hiperboliczna}, zwana też geometrią Łobaczewskiego albo Bolyaia--Łobaczewskiego \cite[s. 287]{eves1_1972}. % lang-pl
Gauss, który w wieku sześćdziesięciu dwóch lat zacznie uczyć się rosyjskiego, prawdopodobnie po to, by czytać pisma z Rosji -- w tym Łobaczewskiego -- doceni go prywatnie. % lang-pl
William Kingdon Clifford nazwie Łobaczewskiego „Kopernikiem geometrii''. % lang-pl
A differenza di Bolyai costruirà soltanto la geometria in cui per un punto passa più di una parallela. % lang-it
Poiché la pubblicherà per primo, sarà il suo nome a essere dato alla \emph{geometria iperbolica}, detta anche geometria di Lobachevskij o di Bolyai--Lobachevskij \cite[p. 287]{eves1_1972}. % lang-it
Gauss, che a sessantadue anni comincerà a studiare il russo, probabilmente per leggere gli scritti provenienti dalla Russia -- tra cui quelli di Lobachevskij -- lo apprezzerà in privato. % lang-it
William Kingdon Clifford chiamerà Lobachevskij il «Copernico della geometria». % lang-it
\index[persons]{Clifford, William Kingdon}%

W 1846 roku Łobaczewski zostanie zwolniony z uniwersytetu z powodu pogarszającego się zdrowia, a w 1856 roku umrze w ubóstwie. % lang-pl
Nel 1846 Lobachevskij sarà licenziato dall'università a causa del peggioramento della salute e nel 1856 morirà in povertà. % lang-it
